%\subsection{Polymorphisme du dioxyde de vanadium}
%
\begin{exercise}[Polymorphisme du dioxyde de vanadium]

L'oxyde de vanadium IV (\ce{V2O4}) existe sous de variétés allotropiques notées
$\alpha$ et $\beta$. Le composé $\ce{V2O4(\beta)}$ est stable au-dessus de \SI{345}{\kelvin}.
La chaleur spécifique molaire du composé $\ce{V2O4(\beta)}$ est supérieure de
\SI{1.25}{\joule\per\mole\per\kelvin} à celle du composé $\ce{V2O4(\alpha)}$ à toute température.

Calculer l’enthalpie libre molaire standard de la transformation ($\Delta_{\alpha \to \beta } G^\circ (298 ~ \si{\kelvin})$ ) :
\begin{align*}
	\ce{V2O4 (\alpha) ->  V2O4 (\beta)}
\end{align*}

Sachant que pour cette transformation $\Delta_{\alpha \to \beta } H^\circ (345 ~ \si{\kelvin}) = 8610 ~ \si{\joule\per\mole}$

\end{exercise}